%% ma: L11-B02-0006
%% lop: 11
%% chuong: 1
%% bai: 2
%% dang: 11.2.9
%% loai: TL
%% muc: VD
%% dap_an: "-"
%% nguon: "de-cuong-giua-ky-1-toan-11-nam-2026-2027-truong-thpt-hoang-van-thu-ha-noi-2-8.pdf (M08, Đề cương giữa kì 1 Toán 11 THPT Hoàng Văn Thụ, Hà Nội), Phần 4 Bài 2 ý 1"
%% kien_thuc: "Bài 1 (giá trị lượng giác của góc lượng giác), Bài 2 (công thức lượng giác)"
%% trang_thai: cho-duyet
%% ly_do: "Dạng mới đề xuất 11.2.9"
%% ghi_chu: "Bổ sung \"giả thiết các biểu thức đều có nghĩa\" (gốc không nêu điều kiện)."
\begin{ex}
Chứng minh rằng (giả thiết các biểu thức đều có nghĩa):

a) $\dfrac{1+\sin2x+\cos2x}{1+\sin2x-\cos2x}=\cot x$;

b) $\dfrac{\cos^2\alpha-\sin^2\alpha}{\cot^2\alpha-\tan^2\alpha}=\sin^2\alpha\cos^2\alpha$;

c) $\dfrac{\tan\alpha-\tan\beta}{\cot\beta-\cot\alpha}=\tan\alpha\tan\beta$;

d) $\dfrac{1-4\sin^2x\cos^2x}{(\sin x+\cos x)^2}=(\sin x-\cos x)^2$.
\loigiai{a) Tử $=2\cos^2x+2\sin x\cos x=2\cos x(\cos x+\sin x)$, mẫu $=2\sin^2x+2\sin x\cos x=2\sin x(\sin x+\cos x)$; thương bằng $\dfrac{\cos x}{\sin x}=\cot x$.

b) $\cot^2\alpha-\tan^2\alpha=\dfrac{\cos^4\alpha-\sin^4\alpha}{\sin^2\alpha\cos^2\alpha}=\dfrac{\cos^2\alpha-\sin^2\alpha}{\sin^2\alpha\cos^2\alpha}$, nên vế trái bằng $\sin^2\alpha\cos^2\alpha$.

c) $\cot\beta-\cot\alpha=\dfrac1{\tan\beta}-\dfrac1{\tan\alpha}=\dfrac{\tan\alpha-\tan\beta}{\tan\alpha\tan\beta}$, nên vế trái bằng $\tan\alpha\tan\beta$.

d) Tử $=1-\sin^22x=(1-\sin2x)(1+\sin2x)$, mẫu $=1+\sin2x$; thương bằng $1-\sin2x=(\sin x-\cos x)^2$.}
\end{ex}
